%=============================================================================
% chktex-file 1 chktex-file 8 chktex-file 12 chktex-file 13 chktex-file 24 chktex-file 26
% RÉSUMÉ & ABSTRACT
%=============================================================================

\chapter*{Résumé}
\addcontentsline{toc}{chapter}{Résumé}

Le secteur du crédit-bail (\textit{leasing}) financier fait face à des défis opérationnels, juridiques et financiers majeurs lors de la gestion du recouvrement contentieux des véhicules en défaut de paiement. Les processus traditionnels, reposant sur des feuilles de calcul locales disparates et des flux documentaires papier, souffrent de délais de traitement excessifs, d'un manque de traçabilité des échéances légales impératives et d'un risque élevé d'erreurs d'évaluation financière lors de la reprise des actifs sur le marché secondaire.

Ce Projet de Fin d'Études, réalisé au sein de \textbf{\hostCompany}, présente la conception et la réalisation de \textbf{\projectShortTitle}, une plateforme logicielle SaaS (\textit{Software as a Service}) B2B multi-tenant dédiée à l'automatisation intégrale du cycle de recouvrement et d'expertise automobile. La solution structure le workflow réglementaire en cinq phases séquentielles (pré-contentieux, mise en demeure, saisie du véhicule, vente et clôture), tout en intégrant un pipeline d'intelligence artificielle asynchrone capable d'extraire automatiquement les données financières des rapports d'expertise PDF non structurés (NLP/LLM) et de confronter instantanément la valeur marchande réelle à la valeur résiduelle contractuelle.

Développée avec une architecture microservices polyglotte moderne associant \textbf{Next.js 16} (BFF sécurisé), \textbf{Spring Boot 4} (Java 21), \textbf{FastAPI} (Python 3.11), \textbf{PostgreSQL 18} avec isolation multi-tenant par schémas étanches et \textbf{Docker}, la plateforme garantit une conformité légale rigoureuse (RGPD / Loi n°63-2004), une traçabilité d'audit immuable et une réduction substantielle des pertes financières pour les institutions de crédit-bail.

\vspace{0.8cm}
\noindent\textbf{Mots-clés :} Crédit-bail automobile, Recouvrement contentieux, SaaS Multi-Tenant, Intelligence Artificielle, LLM, NLP, Next.js, Spring Boot, FastAPI, PostgreSQL.

\newpage

\chapter*{Abstract}
\addcontentsline{toc}{chapter}{Abstract}

The automotive leasing industry faces critical operational, regulatory, and financial challenges in managing debt collection and asset repossession for defaulted contracts. Traditional management workflows, heavily reliant on fragmented spreadsheets, manual paper trails, and unindexed PDF reports, suffer from severe administrative delays, lack of statutory deadline tracking, and significant inaccuracies during asset remarketing.

This graduation thesis, conducted at \textbf{\hostCompany}, introduces \textbf{\projectShortTitle}, an enterprise B2B multi-tenant SaaS (\textit{Software as a Service}) platform designed to digitize and automate the entire vehicle recovery and appraisal lifecycle. The system enforces a strict 5-phase regulatory workflow (pre-litigation, formal notice, vehicle seizure, asset remarketing, and case closure). At its core, an asynchronous AI-powered document intelligence pipeline leverages NLP and Large Language Models (LLMs) to parse unstructured appraisal PDF reports, automatically benchmarking real-time market value against contractual residual value.

Built with a modern polyglot microservices architecture combining \textbf{Next.js 16} (secure BFF layer with HttpOnly cookies), \textbf{Spring Boot 4} (Java 21), \textbf{FastAPI} (Python 3.11), \textbf{PostgreSQL 18} (schema-per-tenant isolation), and \textbf{Docker}, the platform guarantees absolute data confidentiality, immutable audit logging, regulatory compliance (GDPR / Law No. 63-2004), and measurable financial risk reduction.

\vspace{0.8cm}
\noindent\textbf{Keywords :} Automotive Leasing, Debt Recovery, Multi-Tenant SaaS, Artificial Intelligence, LLM, NLP, Document Understanding, Next.js, Spring Boot, FastAPI, PostgreSQL.
